%%=============================================================================
%% Inleiding
%%=============================================================================
\chapter{\IfLanguageName{dutch}{Inleiding}{Introduction}}%
\label{ch:inleiding}
Vesalius.ai ontwikkelt software die het dagelijkse werk van artsen en andere
zorgverleners moet vereenvoudigen. Het platform ondersteunt verschillende
processen binnen de medische praktijk, waaronder de communicatie met patiënten
en de administratieve verwerking van consultaties. Door deze processen te
digitaliseren en te automatiseren, wil Vesalius.ai de administratieve werklast
van zorgverleners verminderen en hen meer tijd laten besteden aan hun kerntaken.
Binnen deze omgeving wordt voortdurend gezocht naar mogelijkheden om de
functionaliteit, configureerbaarheid en gebruiksvriendelijkheid van het
platform verder te verbeteren.

Tijdens mijn stage werkte ik mee binnen het ontwikkelingsteam van Vesalius.ai.
Mijn dagelijkse stagewerk omvatte onder meer het analyseren en oplossen van
bugs en het uitvoeren van andere ontwikkeltaken binnen de bestaande applicatie.
Binnen deze stage werd het graduaatsproject afzonderlijk afgebakend als een
zelfstandig technisch project dat inspeelt op een concrete behoefte van het
bedrijf. De aanleiding voor het project is een bestaand backlog-item rond het
flexibeler beheren van automatisch verstuurde e-mails. Dit onderwerp biedt de
mogelijkheid om zelfstandig een oplossing te analyseren, ontwerpen,
implementeren en testen, zonder dat hiervoor rechtstreeks in de
productieomgeving moet worden gewerkt.

\section{\IfLanguageName{dutch}{Probleemstelling}{Problem statement}}%
\label{sec:probleemstelling}
Binnen het Vesalius-platform worden op vooraf bepaalde momenten automatisch
e-mails verstuurd, bijvoorbeeld bij afspraakbevestigingen, herinneringen en
annulaties. De inhoud van deze e-mails wordt grotendeels door het bedrijf
bepaald. Artsen en organisaties kunnen bepaalde elementen van de huisstijl,
zoals het logo en de kleuren, aanpassen, maar hebben bij de meeste e-mailtypes
onvoldoende mogelijkheden om de tekst en opmaak zelf te wijzigen. Voor
bepaalde conversatiemails bestaat reeds een beperkte vorm van aanpassing via
eigen tekst en placeholders binnen een vaste structuur, maar deze
mogelijkheid is niet beschikbaar voor alle e-mailtypes.

Hierdoor kunnen artsen en organisaties de automatische communicatie niet
volledig afstemmen op hun specifieke behoeften. Een visuele editor kan hiervoor
een oplossing bieden door gebruikers toe te laten de inhoud en opmaak van
e-mails aan te passen zonder dat zij rechtstreeks HTML-code moeten wijzigen.
Een dergelijke oplossing brengt echter ook technische uitdagingen met zich
mee. De gegenereerde HTML moet veilig worden verwerkt, dynamische placeholders
moeten correct worden vervangen en de uiteindelijke e-mails moeten zo
consistent mogelijk worden weergegeven in verschillende e-mailclients, zoals
Gmail, Outlook en Apple Mail.

\section{\IfLanguageName{dutch}{Doelstelling}{Objective}}%
\label{sec:doelstelling}
Het doel van dit graduaatsproject is het ontwikkelen en evalueren van een
zelfstandige proof of concept (PoC) van een WYSIWYG-e-maileditor voor het
aanpassen van automatisch gegenereerde e-mails binnen de context van
Vesalius.ai. De editor moet gebruikers toelaten om de tekst en opmaak van
e-mailtemplates visueel te bewerken, de templates op te slaan en opnieuw te
laden, een voorbeeldweergave te raadplegen en testmails te versturen.

De PoC ondersteunt drie voorbeeld-e-mailtypes: afspraakbevestigingen,
herinneringen en annulaties. Daarbij wordt gebruikgemaakt van fictieve
gegevens en een vaste e-mailschil, waarin onder meer het logo, de kleuren en
de footer behouden blijven. De oplossing wordt gevalideerd aan de hand van
geautomatiseerde tests en controles in gangbare e-mailclients. Het project
is geslaagd wanneer de belangrijkste functionaliteiten aantoonbaar werken,
de HTML-verwerking aan de vooropgestelde veiligheidsvereisten voldoet en de
resultaten van de compatibiliteitstests gedocumenteerd zijn.

Het uiteindelijke resultaat bestaat uit een werkende demonstratie, een
reproduceerbare testopstelling, technische documentatie en een onderbouwd
integratievoorstel. Hiermee krijgt Vesalius.ai inzicht in de technische
haalbaarheid en mogelijke meerwaarde van een visuele e-maileditor als basis
voor een toekomstige uitbreiding van het bestaande platform.

\section{\IfLanguageName{dutch}{Afbakening}{Scope}}%
\label{sec:afbakening}
Het graduaatsproject richt zich uitsluitend op de analyse, het ontwerp, de
realisatie en de validatie van een zelfstandige PoC. De functionele scope
omvat één visuele editor, het aanpassen en opslaan van e-mailtemplates,
ondersteuning voor placeholders, een vaste e-mailschil, HTML-rendering, een
voorbeeldweergave en het versturen van testmails. De implementatie wordt
gedemonstreerd aan de hand van drie voorbeeld-e-mailtypes:
afspraakbevestigingen, herinneringen en annulaties.

De PoC wordt ontwikkeld met fictieve gegevens en een tijdelijke testopstelling.
De frontend en mock-API worden afzonderlijk opgezet, zodat de functionaliteit
kan worden onderzocht zonder afhankelijk te zijn van de productieomgeving.
De opslagoplossing wordt gekozen op basis van de behoeften van de testopstelling
en wordt niet vooraf vastgelegd op basis van de productiedatabase. Bij de
ontwikkeling gaat bijzondere aandacht naar veilige HTML-opschoning,
placeholderverwerking, CSS-inlining en de correcte generatie van zowel een
HTML- als een tekstuele e-mailversie.

De integratie in het bestaande Vesalius-platform valt uitdrukkelijk buiten
de scope. Ook productiegebruik, de effectieve planning en verzending van
e-mails in de productieomgeving, Keycloak-e-mails en een volledig
configureerbare layout-builder worden niet gerealiseerd. De e-mailschil blijft
binnen de PoC grotendeels vast, zodat de focus ligt op het aanpassen van de
inhoud en opmaak van de e-mailbody. De compatibiliteit met verschillende
e-mailclients wordt onderzocht en gedocumenteerd, maar een identieke weergave
in alle clients kan niet op voorhand worden gegarandeerd.

De functionele behoefte komt voort uit de bestaande productcontext en de
besprekingen met het bedrijf. Binnen deze afbakening worden de technische
aanpak, de opbouw van de testopstelling, de implementatie en de validatie
uitgewerkt. De resultaten worden vervolgens gebruikt om aanbevelingen te
formuleren voor een eventuele latere integratie.

\section{\IfLanguageName{dutch}{Stagewerk en graduaatsproject}{Internship work and graduation project}}%
\label{sec:stage-vs-project}
Het graduaatsproject maakt deel uit van mijn stage bij Vesalius.ai, maar beide
hebben een verschillende scope. Tijdens de stage werk ik mee aan uiteenlopende
taken binnen het bestaande ontwikkelingsteam. Deze taken omvatten onder meer
het onderzoeken en oplossen van bugs, het aanpassen van bestaande
functionaliteiten en het leren werken met de codebase, ontwikkelomgeving en
technische werkwijze van het bedrijf.

Het graduaatsproject is hiervan onderscheiden als een zelfstandig afgebakend
project rond een visuele e-maileditor. Waar het dagelijkse stagewerk zich richt
op de noden en prioriteiten van de bestaande applicatie, ligt de nadruk van
het graduaatsproject op het systematisch onderzoeken en realiseren van een
afzonderlijke PoC. Hiervoor worden de vereisten geanalyseerd, technische
keuzes onderbouwd, een oplossing ontworpen en de gerealiseerde functionaliteit
getest en gedocumenteerd.

De oplossing wordt bewust los van de productieomgeving ontwikkeld. Hierdoor
kan de haalbaarheid van de editor, de templateverwerking en de e-mailrendering
worden onderzocht zonder dat de bestaande e-mailinfrastructuur moet worden
aangepast. De zelfstandigheid binnen het graduaatsproject ligt voornamelijk
in de technische uitwerking, de keuze en inrichting van de testopstelling,
de implementatie van de verschillende onderdelen en de systematische
validatie van het eindresultaat. De voortgang en belangrijke ontwerpkeuzes
worden afgestemd met de bedrijfsmentor.

\section{\IfLanguageName{dutch}{Leeswijzer}{Guide to the report}}%
\label{sec:leeswijzer}
Na deze inleiding wordt in hoofdstuk~\ref{ch:technologie} de technologische
context van het project besproken. Daarbij worden de gekozen technologieën,
hun rol binnen de PoC en de belangrijkste technische overwegingen toegelicht.

Hoofdstuk~\ref{ch:ontwerp} behandelt de analyse en het ontwerp van de
oplossing. De functionele en niet-functionele vereisten, de architectuur,
de verwerking van templates en placeholders, de beveiligingsmaatregelen en
de belangrijkste ontwerpkeuzes komen hier aan bod.

In hoofdstuk~\ref{ch:realisatie} wordt de implementatie van de PoC beschreven.
Daarbij wordt ingegaan op de opbouw van de frontend en mock-API, de editor,
de templateverwerking, de renderingpipeline en de inrichting van de
testomgeving.

Hoofdstuk~\ref{ch:resultaat} bespreekt de gerealiseerde functionaliteiten
en de validatie ervan. De testaanpak, de resultaten van geautomatiseerde
tests en de vastgestelde verschillen tussen e-mailclients worden hier
toegelicht. Tot slot bevat hoofdstuk~\ref{ch:reflectie} de conclusie,
de persoonlijke reflectie, de belangrijkste leerervaringen en mogelijke
verdere ontwikkelingen van de oplossing.
